% !TEX root = ../main.tex
\section{Bounding the Request Space}
\label{sec:bounded}

Section~\ref{sec:canonical-form} argued that the absence of a canonical form is what
makes accuracy uncertifiable, and Section~\ref{sec:normaliser} argued that training a
semantic normaliser is not the remedy. There is a third option, and it is the one an
architect reaches for first: if the model will not canonicalise its input and cannot
be taught to, \emph{do not give it free-form input at all}.

Restricting the interface is not a workaround. It is the direct application of the
diagnosis: canonicalisation moves upstream of the model, into a schema, where it is
inspectable and versioned. This section works through what that buys, what it does
not, and where the residual risk goes---because it does not disappear.

\subsection{The construction}

Replace free text with a typed request. Instead of ``is this approved for children?''
the interface accepts an object:

\begin{center}
\small
\begin{tabular}{@{}l>{\raggedright\arraybackslash}p{7.2cm}@{}}
\toprule
\texttt{intent} & \texttt{INDICATION\_CHECK} (from a closed set) \\
\texttt{subject\_id} & a key into a system of record \\
\texttt{indication} & a coded value \\
\texttt{jurisdiction} & required \\
\texttt{as\_of} & required \\
\texttt{population} & required, from a controlled vocabulary \\
\bottomrule
\end{tabular}
\end{center}

The object is rendered into a fixed template by code the organisation controls and
versions. Requests are stateless: one request, one fresh context.

The fields are not arbitrary, and the discipline that selects them is the point rather
than the example. Every component of $c$ appears, because
Section~\ref{sec:resolved-context} established that a claim whose context is unresolved is
not yet evaluable---so jurisdiction, valid time, and intended use are required fields, and
a request missing any of them does not validate. The material facts $f$ enter as keys into
systems of record rather than as free text, so that the entity under discussion is
identified rather than described. The intent is drawn from a closed set, because the
decision class determines which sources have standing under $\standing$ and therefore what
counts as authoritative. And coded values replace prose wherever a controlled vocabulary
exists, since that is what allows applicability to be evaluated as a predicate in
Section~\ref{sec:governed-rag}'s third component rather than judged.

What the schema encodes, in other words, is the set of bindings that would otherwise have
to be guessed. Anything the type system requires is a thing the system can no longer
silently assume.

\subsection{What this fixes}

Against Table~\ref{tab:invariants-system}, the gains are specific and substantial.

\paragraph{I1 becomes satisfied by construction.} Not because the model acquired form
invariance, but because the opportunity to violate it has been removed. Two
materially equivalent requests are now \emph{the same request}: the equivalence class
of Section~\ref{sec:canonical-form} has exactly one member, and the canonical form is
supplied by the schema.

\paragraph{The request space becomes enumerable.} It is a product of typed fields
rather than $\mathcal{V}^{*}$. It may be combinatorially large, but it is
\emph{stratifiable}: coverage can be designed by jurisdiction, date band, population,
and intent, with known sampling properties. That is ordinary test design. The
validation problem becomes well posed in the sense of
Section~\ref{sec:not-easier}---the same sense in which it is well posed for a
classifier.

\paragraph{I10 is enforced by the type system rather than detected.} A request without
a jurisdiction does not validate. The misapplied category of
Section~\ref{sec:cell-v} does not need to be caught downstream, because the binding
it depends on is a precondition of submission. This is the single largest gain, and
it is available with no machine learning at all.

\paragraph{I7 becomes achievable.} Statelessness removes context accumulation, which
is the dominant source of within-session drift. Combined with the reproducibility
conditions of Section~\ref{sec:determinism}, replay becomes a realistic control
rather than an aspiration.

\subsection{What this does not fix}

\paragraph{The output space is still $\mathcal{V}^{*}$.} Bounding the input bounds the
input. The model can still emit any string, so ungrounded assertion survives
untouched: a confident figure resting on nothing is exactly as available as before.
Grounding is I13 and no amount of input discipline supplies it.

\paragraph{Errors are still not label mismatches.} If the response is prose, it can be
partly right, right for the wrong reason, or adequate for a weaker claim than the one
made. The atomic countable error of the classifier case does not follow from a typed
input.

\subsection{Where the problem goes}
\label{sec:relocation}

This is the part that decides whether the design is an improvement or a disguise.
\textbf{Who fills in the request?}

If a person translates their need into the typed object, the ambiguity has not been
eliminated; it has been moved into that translation, and it is unrecorded unless the
system captures it.

If a language model performs the translation---which is what is usually built,
because it preserves the interface that made the technology attractive---then the
unbounded input space has been reintroduced at the parsing stage. Form-invariance
failure now determines \emph{which request object is constructed}. A paraphrase can
produce a different jurisdiction, a different population, or a different intent, and
everything downstream will process that object correctly and produce a well-grounded
answer to the wrong question.

This is worse than the original failure in one specific respect: it is less visible.
An ungrounded assertion at least looks like an assertion, and the audit record shows
what evidence was and was not used. A misparsed request produces a record that is
internally consistent and complete. Every control passes. The evidence applies to the
context in the record, and the context in the record is not the user's context.

The mitigation is not technical. It is to \textbf{render the resolved request back and
require confirmation before proceeding}, which converts a silent failure into a
visible one and puts the binding decision where accountability already sits. This is
the pattern enterprise systems have always used before consequential execution, and
the argument of this paper is that a natural-language front end does not earn an
exemption from it. Where confirmation is impractical, the parse should be treated as
part of the assertion for audit purposes: recorded, versioned, and measured for
accuracy against reviewed samples like any other component.

\subsection{Constraining the output, and what that does and does not buy}
\label{sec:constrained-output}

The natural complement is to constrain the output as well: grammar-constrained
decoding into a schema, with an enumerated verdict field, required evidence
references, and an explicit \texttt{INSUFFICIENT\_EVIDENCE} variant.

This is worth doing and its benefits should be stated exactly, because it is easy to
credit it with more than it delivers.

It \emph{does} guarantee that the output parses, that the output space is enumerable,
and that errors become atomic and countable again---restoring on the output side the
well-posedness that typed requests restore on the input side. It \emph{does} make
abstention a first-class value rather than a hoped-for phrasing, and it makes
per-field evidence linkage (I12) structurally enforceable rather than conventional.

It does \emph{not} satisfy I17. A grammar constrains which strings can be emitted; it
does not determine which the model selects. The choice between a substantive verdict
and the abstention variant is made by the same next-token computation as everything
else, subject to the same pressures identified in Section~\ref{sec:rlhf}. Constrained decoding makes abstention \emph{expressible}, not
\emph{correctly triggered}.

I17 is therefore satisfied by the check, not by the schema. The evidence conditions of
Definition~\ref{def:grounded-evidence} are evaluated outside the model, and the gate
must be able to \emph{override} a substantive verdict the model produced---replacing it
with a qualification, a refusal, or an escalation---when those conditions are not met.
A system that trusts the model's own choice of variant has implemented a vocabulary
for abstention and no gate.

\subsection{Where this is practical}

The design suits high-volume, repetitive, well-typed decisions: eligibility
determinations, coding against a rulebook, standard disclosures, reconciliations,
completeness checks. The user is an operator working a queue against a known decision
class, and the intent set is small and stable.

It suits exploratory work badly, and that matters, because exploration is a large part
of why these systems were adopted. An analyst pursuing an unanticipated question has
no intent to select.

Two costs deserve naming.

\paragraph{The intent taxonomy is an ontology project.} Enumerating decision classes,
coding their fields, and maintaining controlled vocabularies is the same work as
building a conformed dimensional model, with the same effort profile and the same
tendency to be underestimated.

\paragraph{An incomplete taxonomy fails quietly.} A question outside the intent set
does not raise an error. It is routed to the nearest available intent, which is a
misparse of the kind described in Section~\ref{sec:relocation}. Twenty intents may
cover most of the volume, and the residual tail is disproportionately where the
unusual, high-consequence cases live. A closed intent set therefore needs an explicit
\emph{no matching intent} outcome and a route for it, or the constraint will silently
manufacture the failure it was introduced to prevent.

\subsection{The effect on what the system is}

One consequence should be faced directly, because it changes how the component should
be described and governed.

If the input is typed, the output is schema-constrained, and the evidence conditions
are checked outside the model, then it is fair to ask what the model is contributing.
Frequently the honest answer is: retrieval, summarisation, and the rendering of a
determination into readable language---while the determination itself is made by
deterministic rules over structured evidence.

That is a legitimate and valuable architecture, and it is one enterprises already know
how to validate. But it \emph{reclassifies} the system. The model is a retrieval and
presentation layer over a decision service, not the decision-maker, and the assurance
argument should say so. A programme that describes such a system as ``AI-driven
decisioning'' is overstating the model's role and, more importantly, misdirecting its
own validation effort toward the component that is doing the least consequential work.

There is a further consequence worth facing, because it follows from combining this
section with Section~\ref{sec:structured} and a reader will combine them. If structured
evidence is the strongest available form of $E$, and typed requests plus constrained
output are what reach the higher reliance tiers, then the architecture that best
satisfies this paper's conditions is one in which the model does comparatively little of
the deciding. Taken to its conclusion, that is an argument for routing decidable claims
\emph{away} from the model altogether.

I take that to be a feature rather than an embarrassment. A framework that identifies
which decisions do not need a language model is more useful than one that assumes every
decision does. What remains for the model, once the decidable residue is routed away, is
the genuinely open-textured part: synthesis across heterogeneous sources, questions whose
evidence exists only as prose, interpretation where the governing text is ambiguous, and
interaction with people who cannot state their need in a schema. That is a smaller claim
about the technology than the field usually makes, and a more defensible one---and it is
the part where the measurement problem of Section~\ref{sec:auditable} genuinely bites,
because it is the part no schema can decide.

The tiering of Section~\ref{sec:tiers} is where this resolves. Bounding the request
space is not a general answer to the problem of ungrounded assertion; it is the
mechanism by which a narrow, well-understood decision class can be moved to a higher
reliance tier than free-form interaction can support. Free-form interfaces are
appropriate at Tiers~0 to~2. Reaching Tier~3 and above is, in practice, an exercise in
bounding what may be asked and what may be answered.
